\section{Introduction}
\label{sec:intro}

Vision-language models (VLMs) have made remarkable progress in recent years, with systems like GPT-4o~\cite{radford2021clip}, Claude~3.5 Sonnet, and Gemini~\cite{li2023blip2} achieving impressive performance on visual question answering, image captioning, and multimodal reasoning tasks. However, a growing body of evidence suggests that high accuracy on standard benchmarks may not reflect genuine visual understanding~\cite{thrush2022winoground, yuksekgonul2023aro, tong2024eyeswideShut}. Models may exploit linguistic shortcuts, rely on dataset biases, or associate statistical co-occurrences rather than performing true visual reasoning~\cite{hsieh2023sugarcrepe, huang2024conme}.

A fundamental limitation of existing benchmarks is that they evaluate models on individual images in isolation. A model that answers ``the red square is above the blue circle'' may be correct, but we cannot know whether this answer reflects genuine spatial understanding or a lucky guess. Drawing on Pearl's causal hierarchy~\cite{pearl2009causality, pearl2018bookofwhy}, we argue that testing \emph{counterfactual} reasoning---``if the positions were swapped, would the model's answer change?''---provides a more rigorous evaluation of visual understanding. Counterfactual evaluation sits at the highest level of Pearl's ladder of causation, requiring models to reason about what \emph{would} happen under intervention, not merely what \emph{does} happen.

We introduce \method, a counterfactual consistency benchmark for VLMs that operationalizes this insight. The core idea is simple: for each scene, we create a paired \emph{intervened} version with a single, controlled modification (\eg, swapping two objects, removing an arrow, changing a color) and ask the same question about both images. The key test is whether the model's answer changes \emph{if and only if} the ground-truth answer changes. We formalize this as the \emph{Counterfactual Consistency Score} (\ccs), decomposed into \emph{sensitivity} (the model correctly changes its answer when it should) and \emph{specificity} (the model correctly maintains its answer when it should).

\method has several attractive properties:
\begin{itemize}[leftmargin=*,nosep]
    \item \textbf{Cheap.} All 325 scene pairs are generated programmatically with Python's PIL library in under 30 seconds. No annotation, no data collection, no curation.
    \item \textbf{Controllable.} Because we generate both images and ground-truth labels, we have perfect control over what the correct answer is and whether the intervention should change it.
    \item \textbf{Diagnostic.} Five reasoning categories (spatial, causal, compositional, counting, occlusion) and nine intervention types allow fine-grained diagnosis of model weaknesses.
    \item \textbf{Extensible.} New scene types, interventions, and difficulty levels can be added trivially by writing a new Python generator function.
    \item \textbf{Complementary.} \ccs measures something fundamentally different from accuracy: a model can be 100\% accurate on individual images yet have low \ccs if it gives the same answer regardless of intervention.
\end{itemize}

Our experiments on six VLMs---three proprietary and three open-weight---reveal striking variations in counterfactual consistency. \ccs spans a 19.5-point range across models, from 99.7\% (Gemini~2.0 Flash) to 80.2\% (Llama~3.2 11B), even though all models exceed 82\% accuracy on individual images. We uncover two distinct failure profiles: Claude~3.5 Sonnet suffers from \emph{low specificity} (80.7\%), spuriously changing answers for irrelevant interventions, while Llama~3.2 11B suffers from \emph{low sensitivity} (77.4\%), failing to update answers when the visual scene genuinely changes. Notably, the open-weight Qwen2.5-VL~72B achieves 97.5\% \ccs with perfect specificity, rivaling proprietary GPT-4o.

Our contributions are:
\begin{enumerate}[leftmargin=*,nosep]
    \item We formalize the \emph{Counterfactual Consistency Score} (\ccs) and its decomposition into sensitivity and specificity as a principled evaluation metric for VLMs.
    \item We release \method, a fully synthetic, zero-cost benchmark with 325 paired scenes spanning five reasoning categories and nine intervention types.
    \item We evaluate six VLMs and reveal systematic inconsistencies---a 19.5-point \ccs gap and two distinct failure profiles---that standard accuracy metrics cannot detect.
\end{enumerate}
